\section{Compact TT propagation and two-ended closure}
\label{sec:propagation}

Let $q_{ij}$ be a complex STF quadrupole amplitude and $\Lambda(\hat n)$ the transverse-traceless projector for direction $\hat n$.  Define the angular power form
\begin{equation}
F_q(\hat n)=q^*:\Lambda(\hat n):q .
\end{equation}
Since $\Lambda$ is an orthogonal projector, $0\le F_q\le q^*:q$.  For an STF tensor,
\begin{equation}
F_q=q^*:q-2(q^*\hat n)\cdot(q\hat n)
+\frac12|\hat n\cdot q\hat n|^2 .
\end{equation}
Using the standard isotropic second and fourth angular moments gives
\begin{equation}
\boxed{
\int\dd\Omega\,F_q(\hat n)=\frac{8\pi}{5}q^*:q .
}
\label{eq:Fint}
\end{equation}
Hence the normalized angular density
\begin{equation}
w_q(\hat n)=\frac{F_q(\hat n)}{(8\pi/5)q^*:q}
\end{equation}
obeys $w_q\le5/(8\pi)$.  The corresponding directivity is therefore
\begin{equation}
\boxed{D_q=4\pi w_q\le\frac52 .}
\label{eq:Dmax}
\end{equation}
The maximum is attained by a plus quadrupole observed along its symmetry-normal direction.  Gravitational-antenna directivity itself is historical \cite{Hirakawa1976}; Eq.~\eqref{eq:Dmax} is used only to fix the present normalized propagation coefficient.

Choose unit-norm TT polarization tensors $\epsilon^\lambda(\hat n)$ and normalized angular amplitudes
\begin{equation}
u_{q,\lambda}(\hat n)
=\frac{q:\epsilon^\lambda(\hat n)}
{\sqrt{(8\pi/5)q^*:q}}.
\end{equation}
Translation by a separation $R$ adds phase $\exp(ik\hat n\cdot\bm R)$.  In the outgoing wave zone, stationary phase at $\hat n=\hat R$ gives
\begin{equation}
t_{BA}
\sim
\frac{2\pi e^{ikR}}{ikR}
\sum_\lambda
u^*_{B,\lambda}(\hat R)
u_{A,\lambda}(\hat R).
\end{equation}
Cauchy--Schwarz together with $w_q\le5/(8\pi)$ yields
\begin{equation}
\boxed{
\limsup_{kR\rightarrow\infty}
(kR)^2|t_{BA}|^2\le\frac{25}{16} .
}
\end{equation}
Optimizing over normalized compact source and receiver superpositions gives the operator form
\begin{equation}
\boxed{
\limsup_{kR\rightarrow\infty}
(kR)^2\|P_g\|_{\op}^2\le\frac{25}{16} .
}
\label{eq:25over16}
\end{equation}
The same coefficient follows from reciprocal effective area, $D_AD_B(\lambda/4\pi R)^2$, with $D_A,D_B\le5/2$.  An independent canonical one-graviton calculation in the companion repository project reproduces the same leading coefficient and, for an aligned plus-mode specialization, supplies finite-distance corrections; those subleading corrections are not assumed to be universal here.

Combining Eqs.~\eqref{eq:gravcut}, \eqref{eq:narrowresource}, and \eqref{eq:25over16}, with $k_0=\omega_0/c$, gives
\begin{align}
\Gamma_{\rm coh}
&\lesssim
\frac{25}{16(k_0R)^2}
\frac{4G\omega_0^4}{3c^5}
\min(I_{2,A},I_{2,B}) \\
&=
\boxed{
\frac{25G\omega_0^2}{12c^3R^2}
\min(I_{2,A},I_{2,B})} .
\end{align}
This proves Eq.~\eqref{eq:headline} within the declared narrowband leading-wave-zone class.

\subsection{Repeated passive returns}

Let $P_+$ and $P_-$ be forward and reverse separated propagation, and let $R_A,R_B$ be the gravitational reflection blocks of the passive endpoints.  Repeated returns sum to the exact resolvent
\begin{equation}
P_{\rm eff}
=(I-P_+R_AP_-R_B)^{-1}P_+ .
\end{equation}
As subblocks of passive scattering matrices, $\|R_A\|_{\op},\|R_B\|_{\op}\le1$.  With $p_\pm=\|P_\pm\|_{\op}$,
\begin{equation}
\boxed{
\|P_{\rm eff}\|_{\op}
\le\frac{p_+}{1-p_+p_-}
}
\end{equation}
whenever $p_+p_-<1$.  For reciprocal propagation, $p_+=p_-=p$ and $\eta=p^2$, so
\begin{equation}
\|P_{\rm eff}\|_{\op}^2
\le\frac{\eta}{(1-\eta)^2} .
\end{equation}
Since $\eta=O[(kR)^{-2}]$, recurrence changes the \emph{upper ceiling} only at subleading power order and
\begin{equation}
\limsup_{kR\rightarrow\infty}(kR)^2\|P_{\rm eff}\|_{\op}^2\le\frac{25}{16} .
\end{equation}
This is not an equality for the actual recurrent transfer: destructive interference may reduce it.  The scattering composition is established network theory \cite{Redheffer1962}; its role here is only to close a potential passive gravitational loophole.
